\section{Introduction}

Full-waveform inversion reconstructs material properties by repeatedly solving a wave equation and minimizing waveform mismatch~\citep{virieux2009overview}. Its reconstruction quality depends strongly on where sources and receivers are placed, yet measurements and wave simulations are expensive. Optimal experimental design (OED) therefore seeks a small, informative acquisition geometry. In seismic FWI, established methods typically compute sensitivities for a comprehensive candidate survey and greedily select measurements according to spectral or resolution criteria derived from the Gauss--Newton system~\citep{maurer2017oed,krampe2021oed,mercier2025combined}. These approaches are effective, but they require a discrete candidate set, often use greedy selection strategies that can lead to suboptimal designs, and optimize a proxy rather than the actual nonlinear reconstruction outcome.

Differentiable simulators allow a different formulation: treat acquisition
coordinates as continuous variables and optimize them jointly with the FWI workflow. Bilevel learning has been used to select
regularization parameters in inverse problems~\citep{holler2018bilevel} and to
place space--time observations in variational data
assimilation~\citep{castro2020observation}. Related differentiable experimental
design methods use physics-informed or neural
reconstructors~\citep{hemachandra2025pied,darges2025node,liu2025dspo}. Most
closely, \citet{downing2026bilevel} optimizes sensor coordinates for
frequency-domain FWI using supervised bilevel learning, implicit upper-level
gradients, and training velocity models. Our aim is complementary: we study an outer design loop in which the design objective itself is
interchangeable. We compare an oracle design objective, which requires the true velocity model, with practical and less expensive alternatives.


Our application is motivated by guided-wave thickness mapping and scanning laser
Doppler vibrometry (LDV). Previous work has shown that FWI can reconstruct plate thickness from guided-wave measurements~\citep{zuo2022quantitative}, making it relevant to non-destructive testing of thin structures such as plates. LDV enables
highly flexible receiver placement, but a measurement at each spatial position must be repeated many times and stacked to improve its signal-to-noise ratio~\citep{Thomsen2024}.
Acquisition time therefore grows quickly with the number of positions, strongly motivating OED for a reduced acquisition
layout.

This motivates an outer design loop for online acquisition rather than a
single layout designed before the experiment. In practice, the LDV
first measures the wavefield at the current receiver coordinates. A short block of FWI is then performed. The selected design objective uses that partially
updated velocity model to calculate new receiver coordinates. The LDV then moves to
those coordinates, acquires a new stacked data set, and the loop repeats. All
design objectives studied below are embedded in this same outer design loop.
Figure~\ref{fig:diff-oed-workflow} illustrates its
measurement--inversion--design--reacquisition sequence.

% This study presents a proof of concept with three main contributions:
% \begin{itemize}
%   \item a fully differentiable acquisition and acoustic FWI workflow that optimizes continuous receiver coordinates;
%   \item a comparison of an oracle model error design objective with a waveform misfit objective and a stochastic illumination objective in the same computational framework; and
%   \item preliminary evidence that the compared objectives improve reconstruction in a synthetic non-destructive testing setting.
% \end{itemize}

Our contributions are a differentiable acoustic FWI workflow for continuous
receiver placement, a comparison of oracle model error, waveform misfit, and stochastic
illumination objectives, and a synthetic demonstration that each improves
reconstruction.

% More broadly, differentiating the complete acquisition and inversion workflow
% makes its components modular: expensive solver or reconstruction components
% could in principle be replaced by learned surrogates, while internal
% algorithmic hyperparameters could be optimized jointly with the acquisition geometry. This study does not yet explore these possibilities.
